\clearpage
\fancyhead[L]{\small HARDINGS HIMMELSATLAS.}
\hypertarget{native-546-lower}{}\label{native-546-lower}
\begin{center}\small Göttingische gelehrte Anzeigen. Stück 158. S. 1569..1571. 3. Oktober 1812.\\[3pt]\rule{.72\linewidth}{.3pt}\end{center}

Die beiden ersten Lieferungen des großen Himmelsatlas unseres Herrn Professor \textsc{Harding} erschienen in den Jahren 1809 und 1810 und wurden seinerzeit in unseren Blättern besprochen. Wir holen jetzt die Besprechung der dritten, im Jahr 1811 ausgegebenen Lieferung nach und verbinden damit die \SourceLeaf{547}Besprechung der vierten, die soeben erschienen ist. Da wir bei der ersten Lieferung bereits ausführlich über den Plan und die Eigenart dieses Werks berichtet haben, brauchen wir hier nur zu versichern, dass die folgenden Lieferungen ganz im selben Geist und mit derselben Sorgfalt ausgearbeitet sind. An Reichhaltigkeit und Feinheit des Stichs übertreffen sie die erste Lieferung zum Teil sogar noch.

Die \emph{dritte} Lieferung besteht aus den Blättern VII, XVI, XVII und XVIII. Blatt Nr. VII erstreckt sich in der Rektaszension von $239^\circ$ bis $281^\circ$, in der Deklination von $2^\circ$ nordlich bis $32^\circ$ südlich. Es enthält also Teile der Sternbilder Skorpion Schütze, Ophiuchus, Schlange und Sobieskisches Schild. Blatt Nr. XVI hat in der Rektaszension dieselbe Ausdehnung wie Nr. VII und reicht in der Deklination von $1^\circ$ südlich bis $33^\circ$nordlich. Außer dem größten Teil des Herkules enthält es Teile der Schlange, des Ophiuchus, des Poniatovsky'schen Stiers und der Leier. Das außerordentlich reichhaltige Blatt Nr. XVII, zu dem Herr \textsc{von Lindenau} dem Verfasser eine beträchtliche Zahl von Sternpositionsbestimmungen geliefert hat, schließt an das vorherige Blatt an. Es hat dieselbe Ausdehnung in der Deklination und reicht bis $321^\circ$ Rektaszension. Außer dem ganzen Delfin, dem Pfeil und dem Fuchs mit der Gans enthält es den größten Teil des Füllens und des Adlers und Antinous sowie Teile vom Poniatovsky'schen Stier, der Schlange, dem Herkules, der Leier, dem Schwan, Pegasus und Wassermann. Auf gleiche Weise schließt Blatt Nr. XVIII wiederum an das vorherige an und reicht bis zur Rektaszension von $1^\circ$. Darauf finden wir beinahe den ganzen Pegasus sowie Teile vom Schwan, Füllen, Wassermann, den Fischen und der Andromeda.

\newpage
Die \emph{vierte Lieferung} zeigt uns ausschließlich südliche Zonen, für die der Verfasser während seines Aufenthalts an der Mannheimer Sternwarte selbst viele ergänzende Sterne beobachtet hat. Sie besteht aus den Blättern XI, XII, XIII und XIV. Diese Blätter reichen alle von $1^\circ$ nordlich bis $33^\circ$ südlich und schließen so aneinander an, dass sie in der Rektaszension die Zone von $39^\circ$ bis $201^\circ$ umfassen. Wir finden darauf also teils oder ganz den Walfisch, die Georgsharfe, das Laboratorium, den Eridanus, das Brandenburgische Szepter, die Bildhauerwerkstatt, den Orion, den Hasen, die Taube, den Großen Hund, das Schiff Argo, das Einhorn, die Buchdruckerwerkstatt, den Kompass, die Luftpumpe, Hydra, den Sextanten, den Löwen, die Jungfrau, den Becher, den Raben und den Zentaur. Wir bemerken nur noch, dass die \textsc{Vandenhoeck}'sche Buchhandlung von jetzt an den Verschliess der Karten übernommen hat.
